% META: {"id": "TCOMPL-ECH-CO-006", "chapitre": "TCOMPL-ECHANTILLONNAGE", "type_objet": "corrige", "exercice_ref": "TCOMPL-ECH-EX-006", "capacites_codes": ["C2"], "sources_inspiration": [], "status": "generated"}
% BEGIN-VERIFY
% from sympy import binomial
% lignes = [[binomial(n, k) for k in range(n + 1)] for n in range(6)]
% assert lignes[5] == [1, 5, 10, 10, 5, 1]
% assert lignes[4] == [1, 4, 6, 4, 1]
% assert binomial(5, 2) == binomial(5, 3) == 10
% assert binomial(4, 1) + binomial(4, 2) == binomial(5, 2)
% for n in range(1, 12):
%     for k in range(1, n + 1):
%         assert binomial(n, k) == binomial(n - 1, k - 1) + binomial(n - 1, k)
%         assert binomial(n, k) == binomial(n, n - k)
% END-VERIFY
\begin{corrige}{TCOMPL-ECH-EX-006}
\textbf{1.} Les lignes $0$ à $5$ sont $1$ ; $1\ 1$ ; $1\ 2\ 1$ ;
$1\ 3\ 3\ 1$ ; $1\ 4\ 6\ 4\ 1$ ; $1\ 5\ 10\ 10\ 5\ 1$.

\textbf{2.} $\binom52=10$ et $\binom53=10$ : les deux sont égaux, car choisir
$2$ éléments parmi $5$ revient à en écarter $3$. En général
$\binom nk=\binom n{n-k}$.

\textbf{3.} $\binom41+\binom42=4+6=10=\binom52$. C'est la relation de Pascal :
$\binom nk=\binom{n-1}{k-1}+\binom{n-1}{k}$, qui construit chaque ligne à
partir de la précédente.
\end{corrige}
